% titel: Blatt 5, A8: Ionenkristall oder Molekülverbindung?
% skript: Mortimer Kap. 7 (chemie: Ionenbindung, Atombindung)

Ob eine Verbindung als Ionenkristall (=Ionengitter aus Kation und Anion) oder
als Molekül (Atome kovalent gebunden) vorliegt, entscheidet der
Elektronegativitätsunterschied der Bindungspartner und deren Natur.

\begin{werkzeugkasten}[Typ einer Verbindung abschätzen]
\begin{itemize}
  \item \textbf{Ionisch} (Salz): Metall + Nichtmetall, auch
        Ammonium-/Nitrat-Ionen; grosse EN-Differenz $\to$ Elektronen werden
        übertragen, Gitter aus Kation/Anion.
  \item \textbf{Kovalent/Molekül}: Nur Nichtmetalle, kleine EN-Differenz; Atome
        teilen Elektronen. Auch polare Moleküle (z.\ B. $\mathrm{HCl}$) sind keine
        Ionenkristalle, sondern Molekülverbindungen.
  \item Ausnahmen prüfen: Ammoniumnitrat $\mathrm{KNO_3}$ enthält zwei Ionen
        ($\mathrm{K^+}$ und $\mathrm{NO_3^-}$) und ist trotz Nichtmetall-Partner
        ein Salz.
\end{itemize}
\textit{Skript:} Mortimer Kap. 7 (Bindungstypen).
\end{werkzeugkasten}

\begin{achtung}[Darauf achten]
\begin{itemize}
  \item Nicht mit der Elektronegativität \textit{einzelner} Atome argumentieren,
        sondern mit der \textit{Differenz} im Molekül und der Metall/Nichtmetall-
        Zuordnung.
  \item $\mathrm{KNO_3}$ ist ein klassischer Fallstrick: dissoziiert in
        $\mathrm{K^+}+\mathrm{NO_3^-}$ und bildet Ionenkristall.
  \item $\mathrm{HCl}$ ist ein Dipol-Molekül mit stark polarer, aber kovalenter Bindung.
  \item Wasserstoff-Verbindungen sind meist kovalent (Moleküle), auch wenn das
        Partneratom ein Nichtmetall ist.
\end{itemize}
\end{achtung}
